\documentclass[12pt,a4paper]{article}

% ---------------------------------------------------------
% Idioma y tipografía
% ---------------------------------------------------------
\usepackage[utf8]{inputenc}
\usepackage[T1]{fontenc}
\usepackage[spanish]{babel}
\usepackage{mathpazo}
\usepackage{microtype}


% ---------------------------------------------------------
% Matemática y tablas
% ---------------------------------------------------------
\usepackage{amsmath,amssymb,amsthm}
\usepackage{booktabs}
\usepackage{tabularx}
\usepackage{geometry}
\usepackage{enumitem}

\geometry{
	margin=2.5cm
}

\title{Ecuaciones Diferenciales Ordinarias\\Clase 1}
\author{}
\date{}

\begin{document}
	
	\maketitle
	
	
	\section{Introducción}
	
	La física constituye uno de los principales motores para la construcción de ecuaciones diferenciales. Sin embargo, no es la única disciplina en la que aparecen. La química, la biología, la economía y la ingeniería utilizan ecuaciones diferenciales para describir sistemas que evolucionan en el tiempo o en el espacio.
	
	Siempre que una magnitud cambia y nos interesa comprender cómo ocurre esa variación, es muy probable que exista una ecuación diferencial capaz de modelar el fenómeno.
	
	Muchos de ustedes ya han trabajado con ecuaciones diferenciales, aunque probablemente no las hayan identificado como tales. Un ejemplo clásico es el movimiento de una masa unida a un resorte.
	
	\subsection{El oscilador armónico}
	
	Consideremos una masa \(m\) unida a un resorte de constante elástica \(k\).
	
	La segunda ley de Newton establece que
	
	\[
	m\frac{d^2x}{dt^2}=\sum_{i=1}^{N}F_i(x,\dot{x},t).
	\]
	
	En este caso, la única fuerza que actúa sobre el sistema es la fuerza restauradora del resorte, dada por la ley de Hooke:
	
	\[
	F=-kx.
	\]
	Sustituyendo esta expresión en la segunda ley de Newton, obtenemos
	
	\[
	m\frac{d^2x}{dt^2}=-kx.
	\]
	Dividiendo ambos miembros por \(m\),
	
	\[
	\frac{d^2x}{dt^2}=-\frac{k}{m}x.
	\]
	La solución general de esta ecuación puede escribirse como
	
	\[
	x(t)=C_1\cos(\omega t)+C_2\sin(\omega t),
	\]
	donde \(C_1\) y \(C_2\) son constantes arbitrarias y
	
	\[
	\omega=\sqrt{\frac{k}{m}}
	\]
	es la frecuencia angular. Esta solución también puede escribirse utilizando una amplitud y una fase. En efecto, proponemos
	
	\[
	x(t)=A\cos(\omega t+\phi).
	\]
	Aplicando la identidad trigonométrica del coseno de una suma,
	
	\[
	\cos(\alpha+\beta)
	=
	\cos\alpha\cos\beta
	-
	\sin\alpha\sin\beta,
	\]
	obtenemos
	
	\[
	x(t)
	=
	A\cos\phi\cos(\omega t)
	-
	A\sin\phi\sin(\omega t).
	\]
	Comparando esta expresión con
	
	\[
	x(t)=C_1\cos(\omega t)+C_2\sin(\omega t),
	\]
	se concluye que
	
	\[
	C_1=A\cos\phi,
	\qquad
	C_2=-A\sin\phi.
	\]
	Elevando ambas ecuaciones al cuadrado y sumándolas,
	
	\[
	C_1^2+C_2^2
	=
	A^2\left(\cos^2\phi+\sin^2\phi\right)
	=
	A^2.
	\]
	Por lo tanto, tomando \(A\geq 0\),
	
	\[
	A=\sqrt{C_1^2+C_2^2}.
	\]
	
	La fase puede determinarse mediante
	
	\[
	\phi=\operatorname{atan2}(-C_2,C_1),
	\]
	lo que fija correctamente el cuadrante. Así, ambas formas de escribir la solución son equivalentes.
	
	El período de oscilación está dado por
	
	\[
	T=2\pi\sqrt{\frac{m}{k}}.
	\]
	
	Las derivadas de la posición son
	
	\[
	\dot{x}(t)=-A\omega\sin(\omega t+\phi)
	\]
	
	y
	
	\[
	\ddot{x}(t)=-A\omega^2\cos(\omega t+\phi).
	\]
	
	Observando que
	
	\[
	\omega^2=\frac{k}{m},
	\]
	
	se obtiene
	
	\[
	\ddot{x}(t)
	=-\frac{k}{m}A\cos(\omega t+\phi)
	=-\frac{k}{m}x(t),
	\]
	
	lo que demuestra que la función propuesta satisface la ecuación diferencial.
	
	\subsection{Otros ejemplos en física}
	
	La física moderna también está formulada en términos de ecuaciones diferenciales.
	
	Las ecuaciones de Einstein,
	
	\[
	G_{\mu\nu}+\Lambda g_{\mu\nu}
	=
	\frac{8\pi G}{c^4}T_{\mu\nu},
	\]
	
	constituyen un sistema de ecuaciones diferenciales acopladas que describe la geometría del espacio-tiempo.
	
	Por otra parte, la ecuación de Schrödinger para una partícula libre en una dimensión es
	
	\[
	i\hbar\frac{\partial\psi}{\partial t}
	=
	\hat{H}\psi,
	\]
	
	donde el operador hamiltoniano está definido como
	
	\[
	\hat{H}
	=
	-\frac{\hbar^2}{2m}
	\frac{\partial^2}{\partial x^2}.
	\]
	
	Por lo tanto,
	
	\[
	i\hbar
	\frac{\partial\psi(x,t)}{\partial t}
	=
	-\frac{\hbar^2}{2m}
	\frac{\partial^2\psi(x,t)}{\partial x^2}.
	\]
	
	Una solución particular es
	
	\[
	\psi_k(x,t)
	=
	Ae^{i(kx-\omega t)},
	\]
	
	siempre que la frecuencia angular y el número de onda satisfagan la relación de dispersión
	
	\[
	\omega=\frac{\hbar k^2}{2m}.
	\]
	
	\section{Definiciones básicas}
	
	A continuación se introducen algunas definiciones fundamentales que utilizaremos a lo largo del curso.
	
	\subsection{Definición: ecuación diferencial ordinaria}
	
	Una \textbf{ecuación diferencial ordinaria (EDO)} es una ecuación que relaciona una función desconocida de una única variable independiente con una o más de sus derivadas.
	
	En forma general,
	
	\[
	F\left(
	x,
	y,
	\frac{dy}{dx},
	\frac{d^2y}{dx^2},
	\dots,
	\frac{d^m y}{dx^m}
	\right)
	=
	0.
	\]
	
	Aquí,
	
	\[
	y=y(x)
	\]
	
	es la función desconocida y \(x\) es la variable independiente.
	
	\subsection{Definición: orden de una ecuación diferencial}
	
	El \textbf{orden} de una ecuación diferencial es el orden de la derivada de mayor orden que aparece en la ecuación.
	
	Por ejemplo,
	
	\[
	\frac{dy}{dx}+y=0
	\]
	
	es una ecuación diferencial de primer orden, mientras que
	
	\[
	\frac{d^2y}{dx^2}+y=0
	\]
	
	es una ecuación diferencial de segundo orden.
	
	\subsection{Definición: forma normal}
	
	Una ecuación diferencial de orden \(m\) está en \textbf{forma normal} cuando la derivada de mayor orden aparece despejada.
	
	Es decir,
	
	\[
	\frac{d^m y}{dx^m}
	=
	f\left(
	x,
	y,
	\frac{dy}{dx},
	\dots,
	\frac{d^{m-1}y}{dx^{m-1}}
	\right).
	\]
	
	En particular, para una ecuación de primer orden,
	
	\[
	\frac{dy}{dx}
	=
	f(x,y).
	\]
	
	\subsection{Definición: solución de una ecuación diferencial}
	
	Una función
	
	\[
	y=\varphi(x)
	\]
	
	es una \textbf{solución} de una ecuación diferencial en un intervalo \(I\) si:
	
	\begin{enumerate}[label=\arabic*.]
		\item posee todas las derivadas requeridas por la ecuación en \(I\);
		\item al sustituirse en la ecuación diferencial, la satisface para todo \(x\in I\).
	\end{enumerate}
	
	\subsection{Definición: intervalo de definición}
	
	Se denomina \textbf{intervalo de definición} al intervalo en el cual la solución y las funciones que intervienen en la ecuación diferencial están definidas y poseen la regularidad necesaria.
	
	\subsection{Definición: solución general}
	
	En muchos casos regulares, la \textbf{solución general} de una ecuación diferencial de orden \(m\) puede expresarse localmente como una familia de soluciones que depende de \(m\) constantes arbitrarias e independientes.
	
	\subsection{Definición: solución particular}
	
	Una \textbf{solución particular} se obtiene a partir de la solución general al asignar valores específicos a sus constantes arbitrarias.
	
	\section{Problemas de valor inicial}
	
	Un \textbf{problema de valor inicial (PVI)} consiste en una ecuación diferencial acompañada por condiciones impuestas sobre la solución y sus derivadas en un mismo punto.
	
	Para una ecuación diferencial de primer orden,
	
	\[
	\frac{dy}{dx}
	=
	f(x,y),
	\]
	
	una condición inicial tiene la forma
	
	\[
	y(x_0)=y_0.
	\]
	
	Por lo tanto, el problema de valor inicial se escribe como
	
	\[
	\begin{cases}
		\dfrac{dy}{dx}=f(x,y),\\[0.3cm]
		y(x_0)=y_0.
	\end{cases}
	\]
	
	Para una ecuación diferencial de orden \(m\),
	
	\[
	\frac{d^m y}{dx^m}
	=
	f\left(
	x,
	y,
	\frac{dy}{dx},
	\dots,
	\frac{d^{m-1}y}{dx^{m-1}}
	\right),
	\]
	
	un problema de valor inicial queda determinado al especificar
	
	\[
	y(x_0)=y_0,
	\]
	
	\[
	y'(x_0)=y_1,
	\]
	
	\[
	y''(x_0)=y_2,
	\]
	
	\[
	\vdots
	\]
	
	\[
	y^{(m-1)}(x_0)=y_{m-1}.
	\]
	
	Es decir, se especifican el valor de la función y los valores de sus primeras \(m-1\) derivadas en un mismo punto \(x_0\).
	
	En conjunto,
	
	\[
	\begin{cases}
		\dfrac{d^m y}{dx^m}
		=
		f\left(
		x,
		y,
		y',
		\dots,
		y^{(m-1)}
		\right),\\[0.3cm]
		y(x_0)=y_0,\\
		y'(x_0)=y_1,\\
		\hspace{0.4cm}\vdots\\
		y^{(m-1)}(x_0)=y_{m-1}.
	\end{cases}
	\]
	
	Estas condiciones permiten, cuando existe unicidad, seleccionar una única solución dentro de la familia de soluciones de la ecuación diferencial.
	
	\section{Ecuaciones diferenciales lineales}
	
	Una ecuación diferencial de orden \(m\) es \textbf{lineal} si puede escribirse en la forma
	
	\[
	a_m(x)\frac{d^m y}{dx^m}
	+
	a_{m-1}(x)\frac{d^{m-1}y}{dx^{m-1}}
	+
	\cdots
	+
	a_1(x)\frac{dy}{dx}
	+
	a_0(x)y
	=
	g(x),
	\]
	
	donde las funciones
	
	\[
	a_0(x),a_1(x),\dots,a_m(x)
	\]
	
	y \(g(x)\) son funciones conocidas, y
	
	\[
	a_m(x)\neq 0
	\]
	
	en el intervalo de interés.
	
	La función desconocida \(y\) y sus derivadas aparecen únicamente de manera lineal: no aparecen productos entre ellas, potencias distintas de uno ni funciones no lineales de \(y\) o de sus derivadas.
	
	Por ejemplo,
	
	\[
	y''+x y'-3y=\sin x
	\]
	
	es lineal, mientras que
	
	\[
	y''+y^2=0
	\]
	
	no es lineal.
	
	\subsection{Definición: ecuación lineal homogénea y no homogénea}
	
	Una ecuación diferencial lineal
	
	\[
	a_m(x)y^{(m)}
	+
	a_{m-1}(x)y^{(m-1)}
	+
	\cdots
	+
	a_1(x)y'
	+
	a_0(x)y
	=
	g(x)
	\]
	
	se denomina \textbf{homogénea} si
	
	\[
	g(x)=0,
	\]
	
	es decir,
	
	\[
	a_m(x)y^{(m)}
	+
	a_{m-1}(x)y^{(m-1)}
	+
	\cdots
	+
	a_0(x)y
	=
	0.
	\]
	
	Se denomina \textbf{no homogénea} si
	
	\[
	g(x)\neq 0.
	\]


		
\end{document}
